\section{Gerbang Sirkuit Digital: Logika pada Keping Silikon}

Pertanyaan mendasar adalah bagaimana Aljabar Boolean yang dipahami melalui konsep dalam buku Boole dapat diimplementasikan dalam perangkat keras komputer.

Jawabannya berkaitan dengan perkembangan penting arsitektur fisika pada abad ke-20, yaitu penemuan \textbf{transistor} semikonduktor yang menerjemahkan nilai abstrak Boolean 1 (True) dan 0 (False) menjadi tegangan listrik, yaitu tegangan tinggi (misalnya 5 Volt) dan tegangan rendah (0 Volt) \cite{grinnell_logic,hackatronic}.

\subsection{Dasar Gerbang Logika (\textit{Logic Gates})}

Transistor berukuran nano disusun membentuk \textit{Logic Gates} yang mengimplementasikan operasi proposisional dasar AND, OR, dan NOT:

\begin{itemize}
    \item \textbf{Gerbang AND (Konjungsi):} Dua input (A, B). Output bernilai 1 hanya apabila kedua input bernilai 1. Jika salah satu input bernilai 0, output bernilai 0.
    \item \textbf{Gerbang OR (Disjungsi):} Dua input (A, B). Output bernilai 1 apabila minimal satu input bernilai 1.
    \item \textbf{Gerbang NOT (Inverter):} Satu input. Output membalik nilai input: input 1 menghasilkan output 0, input 0 menghasilkan output 1.
\end{itemize}

\begin{figure}[H]
\centering
\includegraphics[width=0.95\textwidth]{logic-gates}
\caption{Simbol standar gerbang logika: NOT (kiri), AND (tengah), OR (kanan) dalam arsitektur CPU \cite{hackatronic}.}
\label{fig:logic-gates}
\end{figure}

\subsection{Optimasi Rangkaian: Efisiensi Biaya melalui Aljabar Boolean}

Mengapa perancang perangkat keras perlu menguasai Aljabar Boolean dalam perancangan rangkaian digital?

Diberikan spesifikasi sirkuit: \textbf{$Y = (A \land B) \lor (A \land \neg B)$}. Tanpa analisis, rancangan awal memerlukan 2 Gerbang AND, 1 Gerbang NOT, dan 1 Gerbang OR---total 4 komponen \cite{levin_discrete}.

Insinyur yang menguasai Aljabar Boolean melakukan penyederhanaan:
$$ Y = A \land (B \lor \neg B) \equiv A \land \text{True} \equiv \textbf{A} $$

Ekspresi dapat disederhanakan menjadi A saja. Rangkaian hanya memerlukan satu kabel penghubung dari input A ke output Y, tanpa gerbang AND, OR, maupun NOT. Biaya material dapat dikurangi secara signifikan.
Ini merupakan penerapan Aljabar Proposisional: pengurangan \textbf{BOM (Bill of Materials)} melalui penyederhanaan ekspresi Boolean \cite{grinnell_logic,hackatronic}.
